\chapter{Initial conditions, soundings and prescribed environments}\label{initialization-theory}

An initial condition specifies the experiment from which the equations evolve.
It includes atmospheric thermodynamics, wind, geometry and any deliberate
trigger. Those choices determine the initial inventories and available energy;
they are not neutral preparation for a later physics calculation. Nimbus
separates a sounding source, a hodograph source and a trigger so that choosing
microphysics does not secretly choose an unrelated environment or disturbance.

\section{Three independent questions precede the first timestep}

The sounding source supplies a column of pressure, temperature/theta, moisture
and possibly wind. A file reader interprets a supplied profile; a parametric
source constructs an idealized one. The hodograph policy decides whether to
retain available file winds or build a separate environmental wind. Finally,
the trigger applies no disturbance, a warm bubble or a vortex seed. These are
separate mathematical decisions even when a preset chooses all three.

\begin{figure}[htbp]
\centering
\begin{tikzpicture}[
 box/.style={draw=NimbusBlue,fill=NimbusLight,rounded corners,align=center,
             text width=3.3cm,minimum height=14mm,font=\small},
 flow/.style={-{Latex[length=2mm]},draw=NimbusBlue,thick},
 node distance=14mm and 8mm]
 \node[box] (column) {Sounding source\\file or parameters};
 \node[box,right=of column] (wind) {Hodograph\\file, anchors or zero};
 \node[box,right=of wind] (geometry) {Physical geometry\\heights, faces and terrain};
 \node[box,below=of wind] (placed) {Place model state\\units and EOS};
 \node[box,below=of placed] (trigger) {Selected trigger\\none, bubble\\or vortex};
 \node[box,below=of trigger] (initial) {Initial state\\and inventories\\closure and validation};
 \draw[flow] (column.south)|-(placed.west);
 \draw[flow] (wind)--(placed);
 \draw[flow] (geometry.south)|-(placed.east);
 \draw[flow] (placed)--(trigger);
 \draw[flow] (trigger)--(initial);
\end{tikzpicture}
\caption{Logical responsibilities in constructing an experiment. Environmental
thermodynamics, wind and geometry are separate inputs. A prescribed trigger
can change the initial water and energy inventories; it is not automatically
a conservative microphysical transfer. Exact execution and admission belong
to the initialization and configuration owners.}
\end{figure}


The \href{../include/init/sounding/sounding_source.hpp}{sounding contract},
\href{../include/init/hodograph/hodograph_source.hpp}{hodograph contract} and
\href{../include/init/trigger/trigger_source.hpp}{trigger contract} provide the
interfaces. \href{../src/init/scheme_profile.cpp}{Profile selection} and the
configuration compiler connect them to the ordinary runtime. A helper's
existence does not imply that every option it accepts is available through an
admitted public profile.

An initial environment also needs physical coordinate locations. A terrain
column may begin above the datum of an observed sounding, and its scalar and
velocity levels need not coincide. The initialized pressure, density and theta
must be interpreted through the production dry rho--theta relation. Moist
parcel diagnostics may use virtual temperature while the dynamical EOS uses a
dry pressure law; these are distinct approximations whose roles must be stated.

\section{Reading a sounding is a unit and representation conversion}

The public reader catalog offers \texttt{auto}, \texttt{cm1} and
\texttt{sharpy}. Automatic selection examines content and file signatures.
The reader family includes text and supported netCDF/HDF5 profile forms; a
filename extension alone does not identify a valid physical sounding.
The \href{../src/soundings}{sounding module} validates the supplied variables,
finite values, ordering, units, dimensions and model coverage before a profile
is used.

The mathematical input is a set of values at stated heights or pressure levels,
not a picture of a skew-T diagram. For example, a value written in grams of
vapor per kilogram of dry air must be divided by 1000 to become the solver's
kg/kg dry-air mixing ratio. A temperature in degrees Celsius must be converted
to kelvin before evaluating an exponential saturation formula. Confusing a
moist-air specific humidity $q_s=m_v/(m_d+m_v)$ with a dry-air mixing ratio
$r_v=m_v/m_d$ gives the wrong quantity even though both are dimensionless:
\[
 r_v=\frac{q_s}{1-q_s},\qquad q_s=\frac{r_v}{1+r_v}.
\]
A reader must know which definition the source supplies rather than infer it
from a small decimal value.

Strict-v1 file initialization retains automatic format routing, linear
interpolation, no extrapolation and no silent procedural fallback. The
lower-level fallback source remains compatibility code: its existence does
not authorize replacing a missing observation with an invented atmosphere
under the ordinary strict path. Sounding coverage of the required model heights
is therefore a real input requirement.

For a scalar quantity $a$ between valid source levels $z_j$ and $z_{j+1}$,
linear interpolation uses
\[
 a(z)=(1-\lambda)a_j+\lambda a_{j+1},\qquad
 \lambda=\frac{z-z_j}{z_{j+1}-z_j}.
\]
This preserves the endpoint values and is bounded between them for
$0\le\lambda\le1$. It is not a conservative remapping of a three-dimensional
mass inventory. That is acceptable only because it is being used to define a
new initial profile, with its resulting inventories diagnosed afterward.
Interpolating Cartesian wind components directly also avoids interpolating
across the discontinuity between directions near $0$ and $360$ degrees.

\section{Hydrostatic preparation and the model EOS must be distinguished}

For an idealized resting column, hydrostatic balance is
$dp/dz=-\rho g$. Combining an ideal-gas relation with a chosen temperature
profile gives an integral pressure relation, schematically
\[
 \log\frac{p(z_2)}{p(z_1)}
 =-\int_{z_1}^{z_2}\frac{g}{R_dT_v(z)}\,dz
\]
when the chosen preparation uses virtual temperature $T_v$. The dry version
uses $T$ instead. This equation states the assumptions used to build a profile;
it does not replace the solver's declared dry-density pressure law or prove
that a sampled profile is a discrete resting state.

Nimbus's parametric source constructs pressure along its idealized column and
diagnoses thermodynamic quantities from that construction. File input has its
own conversion and validation. Placement onto model fields, reference-state
construction and later EOS projection must remain consistent with their
respective definitions. A continuous hydrostatic relation and the discrete
pressure-gradient/gravity pairing are separate checks, as developed in the grid
chapter. A small initialization adjustment wave can reflect a mismatch between
them rather than an instability of the supplied atmosphere.

\section{CAPE targeting is a bounded inverse problem}

The \texttt{parametric\_cape} source builds an idealized mixed layer, a
free-tropospheric temperature lapse and a stable upper cap. Its moisture profile
uses the configured near-surface amount, a vertical decay and a relative-humidity
limit. The control variable used to fit CAPE is the free-tropospheric lapse rate,
not an arbitrary independent adjustment of every level.

A surface-parcel buoyancy diagnostic has the conceptual form
\[
 B(z)=g\frac{T_{v,p}(z)-T_{v,e}(z)}{T_{v,e}(z)},\qquad
 \mathrm{CAPE}=\int_{\rm LFC}^{\rm EL}\max(B,0)\,dz.
\]
The parcel and environment must use the same pressure and humidity conventions.
Nimbus's diagnostic estimates the lifting-condensation level, lifts the parcel
dry below it and with a retained moist lapse approximation above it, and
integrates resolved positive buoyancy. It also reports inhibition and the
associated characteristic heights. These quantities depend on the diagnostic
algorithm and vertical resolution, not just the name of the sounding.

The source bisects a lapse-rate interval from 0.0055 to
$0.0098\,\mathrm{K\,m^{-1}}$, seeking agreement within
$50\,\mathrm{J\,kg^{-1}}$, with at most 60 iterations. These are the current
implementation's search bounds and stopping rule. The final profile is still
returned after the iteration limit; the source alone is not a proof that every
requested target is realizable within that family. Requested and realized
CAPE must be distinguished and the applicable diagnostics/admission checked.
See \href{../src/init/sounding/sources/parametric_cape.cpp}{the CAPE source}
and \href{../src/init/sounding/diagnostics.cpp}{parcel diagnostics}.

The \texttt{parametric\_targets} branch translates a requested diagnostic set
into parameters for that same family. Surface moisture may be inferred from an
LCL target; a mixed-layer cap provides an inhibition control; requested LFC/EL
values influence the mixed-layer and upper-profile construction. Those controls
are coupled. CAPE, CIN, LCL, LFC and EL cannot generally be prescribed as five
independent exact numbers by this limited profile family. The
\href{../src/init/sounding/sources/parametric_targets.cpp}{target translator}
contains heuristics and explicit input bounds; the resulting diagnostics remain
the report of what was actually constructed.

\href{https://journals.ametsoc.org/abstract/journals/mwre/110/6/1520-0493_1982_110_0504_tdonsc_2_0_co_2.xml}{Weisman
and Klemp (1982)} supplies important context for studying buoyancy and shear in
idealized convective environments. A preset inspired by that literature is not
a claim that Nimbus reproduces the paper's entire environmental construction,
numerical model or storm evolution.

\section{Wind profiles and hodograph precedence}

A hodograph is the path traced by $(u(z),v(z))$ with height. Its differences
control vertical shear; its curvature affects storm-relative flow and rotation
diagnostics. Translating every wind vector by a constant changes storm-relative
interpretation unless the storm-motion reference is translated consistently.
It does not change the environmental vector shear between two levels.

Nimbus's \texttt{auto} hodograph policy uses sounding winds when they are
available. Otherwise it constructs the configured parametric profile.
\texttt{wk\_param} explicitly selects the parameterized wind, and
\texttt{zero} supplies a calm column. A file containing valid winds and an
explicit separate hodograph are therefore different experiments from the
automatic file-wind path.

The retained \texttt{wk\_param} implementation is piecewise linear in each
Cartesian component through anchors at the surface, 1 km and 6 km. Above 6 km
it holds the final wind constant. Between anchors $a,b$ it uses
\[
 (u,v)(z)=(1-\lambda)(u,v)_a+\lambda(u,v)_b,
 \qquad\lambda=\frac{z-z_a}{z_b-z_a}.
\]
The selector name is not a promise of an arbitrary smooth quarter-circle or a
complete reproduction of a published hodograph. The
\href{../src/init/hodograph/sources/wk_param.cpp}{implemented interpolation}
and \href{../src/init/hodograph/factory.cpp}{wind precedence} are explicit.
Sampling the environmental wind at the chosen scalar/face positions belongs to
initialization; it must not silently change the staggering contract used by
later transport.

\section{The warm bubble changes both heat and, conditionally, water}

The current Cartesian bubble helper uses distance $r$ from its specified
physical center and a compactly applied Gaussian factor
\[
 f(r)=\begin{cases}
       \exp(-9r^2/R^2),&r\le R,\\
       0,&r>R,
      \end{cases}
 \qquad\Delta\theta=\Delta\theta_{\max}f(r).
\]
Thus $R/3$ is the Gaussian width in this convention, and the field is cut off
at $R$ rather than extending to infinity. The helper retains a minimum radius
of 100 m and uses physical terrain height where active metrics are present.
Public validation and exact-profile admission additionally govern which settings
can reach that helper.

When pressure and vapor fields have matching shape and the applied theta
increment is positive, the helper scales vapor to retain its local saturation
ratio at the held pressure:
\[
 q_v^{\rm new}=q_v^{\rm old}
       \frac{q_{vs}(T^{\rm new},p)}{q_{vs}(T^{\rm old},p)}.
\]
The ratio is applied only when the old saturation amount exceeds the retained
small-value threshold. The helper uses its local liquid-water saturation fit
and an Exner exponent of 0.2857 for this conversion. Those are its actual
approximations; one should not assume that it calls the selected microphysics
scheme's saturation routine or performs a closed moist adjustment.

Optional cloud seeding separately adds a configured cloud-water amount times
$f(r)$. It is not subtracted from vapor as a condensation transaction. The warm
bubble therefore changes the initial thermal state and can add vapor and cloud
water. These are prescribed initial-condition changes. Their resulting mass and
energy belong in the initial inventory, rather than being mislabeled as a
conservative source during the first timestep. They also do not guarantee that
every chosen environment will produce sustained convection.

See \href{../src/core/orchestration/dynamics/initial_conditions_cartesian.cpp}{the
Cartesian initializer}. A trigger named \texttt{none} makes no such disturbance;
choosing it does not remove environmental instability or subsequent stochastic
forcing selected elsewhere.

\section{The Rankine seed has a defined profile and an implementation gap}

The \texttt{vortex\_seed} helper adds an azimuthal velocity around a horizontal
center. With radius $r_m$ of maximum wind and signed peak speed $V_m$,
\[
 v_\vartheta(r)=\begin{cases}
 V_m r/r_m,&0<r\le r_m,\\
 V_m r_m/r,&r>r_m,
 \end{cases}
 \qquad
 \Delta u=-v_\vartheta\frac{y-y_c}{r},\quad
 \Delta v=v_\vartheta\frac{x-x_c}{r}.
\]
The axis receives no singular increment. The source adds this wind to the
existing environmental field and applies a hard upper-height cutoff at
\texttt{z\_top\_m}. It does not construct an accompanying balanced pressure
perturbation. Sampling a continuous rotational formula also does not, by itself,
prove zero discrete divergence or exact balance on the selected staggered grid.

The public configuration currently parses and serializes
\texttt{center\_z\_m}, but this helper does not use it to place the vortex
vertically. Its applied placement uses the horizontal center and the upper
cutoff described above. This is a known initialization implementation gap;
the parameter must not be interpreted as an implemented vertical translation.
Exact admitted profiles still have their own restrictions, so this description
is not blanket admission of every vortex configuration. The concrete owner is
\href{../src/init/trigger/sources/vortex_seed.cpp}{the vortex-seed implementation}.

\section{Terrain shapes are also initial scientific inputs}

The terrain catalog provides \texttt{none}, \texttt{bell} and \texttt{schar}.
The flat choice does not imply uniformly spaced vertical levels; the geometry
owner can still provide stretched physical coordinates. A two-dimensional bell
ridge uses
\[
 h(x)=\frac{h_0}{1+(x/a)^2},
\]
while the axisymmetric bell uses
$h(x,y)=h_0a^3/(a^2+x^2+y^2)^{3/2}$. The Sch\"ar profile is
\[
 h(x)=H\exp[-(x/a)^2]\cos^2(\pi x/\ell).
\]
Height and width parameters have physical length units, and the source supplies
the corresponding slopes. These shapes define geometric input, not a physical
land-surface evolution model.

The mapping must turn that topography and the vertical grid into positive
physical layer volumes and compatible face metrics. A wind tangent to a slope,
an environmental profile sampled at the actual height, a surface flux per true
area and a radiation shadow are all consequences of the same terrain geometry.
They must not each reconstruct a different terrain from nominal spacing.
See \href{../src/terrain/base/topography.cpp}{the shape functions} and the
\hyperref[grid-deep-dive]{detailed grid chapter}.

\section{What should be recorded about an initialized experiment}

Retain the source profile or parameter values, wind-selection precedence,
physical grid and terrain, trigger settings, stochastic seed/member identity,
requested diagnostics and realized diagnostics. Record the initial dry mass,
water and energy basis after all prescribed interventions and closures. The
same file profile interpreted on a different grid, with a different trigger or
a different moisture convention is a different experiment even when its name
is unchanged.

The initial-condition and sounding tests under
\href{../tests/init}{tests/init} exercise parsing, thermodynamic conversion,
coverage, selection and trigger independence. The first evolution step then
checks a different contract: it must advance that accepted state through the
selected numerical and physical owners. Restart is different again; it restores
a continuation state and should not rerun an initial bubble or invent a new
sounding in place of saved scheme-private history.
